\section{Công thức toán}

% --- INLINE MATH ---
% Dùng \(...\) hoặc $...$ cho công thức trong dòng
Inline: \( f(x) = x^2 + 2x + 1 \)\\
\indent Hoặc: $f(x) = x^2 + 2x + 1$

% --- DISPLAY MATH (không số) ---
% Dùng \[...\] cho công thức nằm riêng, không đánh số
\[
    \text{ReLU}(x) = \max(0, x)
\]

% --- HỆ PHƯƠNG TRÌNH DÙNG ALIGNED ---
% equation + aligned để vừa căn chỉnh vừa có số thứ tự
\begin{equation}
    \begin{aligned}
        a_1x_1 + a_2x_2 + \dots + a_nx_n & = u \\
        b_1x_1 + b_2x_2 + \dots + b_nx_n & = v \\
        c_1x_1 + c_2x_2 + \dots + c_nx_n & = w
    \end{aligned}
\end{equation}

% --- PHƯƠNG TRÌNH DÀI ---
% dmath tự động xuống dòng (cần gói breqn)
\begin{dmath}
    y = a_1x_1 + a_2x_2 + a_3x_3 + a_4x_4 + a_5x_5 + a_6x_6
      + a_7x_7 + a_8x_8 + a_9x_9 + a_{10}x_{10} + a_{11}x_{11}
      + a_{12}x_{12} + a_{13}x_{13} + a_{14}x_{14} + a_{15}x_{15}
\end{dmath}

% --- HÀM NHIỀU ĐIỀU KIỆN ---
\begin{equation}
    f(x) =
    \begin{cases}
        x^2, & x \ge 0 \\
        -x,  & x < 0
    \end{cases}
\end{equation}

% --- NHIỀU PHƯƠNG TRÌNH ĐỘC LẬP ---
% align giúp căn dấu "=" thẳng hàng
\begin{align}
    a & = b + c \\
    x & = y + z
\end{align}

% --- HỆ PHƯƠNG TRÌNH DẠNG NGOẶC NHỌN ---
\begin{equation}
    \left\{
    \begin{aligned}
        a_1x + b_1y & = c_1 \\
        a_2x + b_2y & = c_2
    \end{aligned}
    \right.
\end{equation}
